%%%%%%%%%%  scan idx 181  |  folio 168  %%%%%%%%%%
\[
\mathrm D^{-}(\Coh(S))\longrightarrow \mathrm D(\Qcoh(S))
\]
\emph{est pleinement fidèle et a pour image essentielle} $\mathrm D^{-}(\Qcoh(S)_{\coh}$.

\emph{Preuve.} On applique (1.1) en prenant pour $S$ le site zariskien de $S$, pour $\mathcal C$ (resp. $\mathcal C_0$) la catégorie fibrée $U\mapsto\Qcoh(U)$ (resp. $\Coh(U)$). Comme $\mathcal O_S$ est cohérent, les objets de $\mathcal C_0$ sont pseudo-cohérents (I 3.5) donc un objet de $\mathrm D(\Qcoh(U))$ ($U$ = ouvert de $S$) est pseudo-cohérent (au sens (2.2.1)) si et seulement s'il l'est relativement à $\mathcal C_0$ (I 2.14.1). Il suffit donc en vertu de (1.1 c)) de montrer que les conditions (i) et (ii) non respée de (1.1) sont satisfaites. Pour (ii), c'est trivial, puisqu'un Module quasi-cohérent qui est quotient d'un Module cohérent est cohérent, et que la notion de cohérence est locale. Pour (i), on utilise (2.1.1 a)) et le fait que tous $\mathcal O_S$-Module quasi-cohérent est limite inductive filtrante de $\mathcal O_S$-Modules cohérents (EGA I 9.4.9).

Admettant (3.5) et (3.7), indépendants du présent numéro, on en déduit facilement :

\textbf{Corollaire 2.2.2.1.} \emph{Soit $S$ un schéma noethérien (resp. noethérien et de dimension finie). Le foncteur canonique}
\[
\mathrm D^b(\Coh(S))\longrightarrow\mathrm D(S)
\]
\[
(\text{resp. }\mathrm D^{-}(\Coh(S))\longrightarrow\mathrm D(S))
\]
\emph{est pleinement fidèle et a pour image essentielle} $\mathrm D^b(S)_{\coh}$ (resp. $\mathrm D^{-}(S)_{\coh}$) (I 2.15).

\textbf{Proposition 2.2.3} (cf. EGA II 4.5.2 et 4.5.5). \emph{Soit $S$ un schéma quasi-compact et quasi-séparé, et soit $(L_i)_{i\in I}$ une famille de faisceaux inversibles sur $S$. Pour tout $\mathcal O_S$-Module $F$, et tout couple $(i,n)\in I\times\mathbb Z$, on note $F_{(i,n)}$ le sous-$\mathcal O_S$-Module de $F\otimes L_i^{\otimes n}$ engendré par les sections de $F\otimes L_i^{\otimes n}$ au-dessus de $S$. Les conditions suivantes sont équivalentes} :
\begin{enumerate}
